\section{Statistical Framework for Quantum Circuit Simulations}
\label{appendix:stats-analysis}

The results presented in Section~\ref{sec:qiskit_sim} of the main paper are derived from simulations of quantum circuits executed on Qiskit's `aer\_simulator` and on real quantum hardware. Each execution, or "shot," is a probabilistic measurement of the system's $(N+1)$ qubits in the computational basis. This appendix details the statistical methodology used to calculate the expectation values of observables and their associated uncertainties.

\subsection{Core Assumptions}
Our analysis is based on the following standard assumptions for quantum circuit simulations:
\begin{itemize}
    \item \textbf{Independent and Identically Distributed (i.i.d.) Trials:} Each shot is an independent measurement drawn from the same underlying probability distribution defined by the final quantum state. The outcome of one shot does not influence any other.
    \item \textbf{Multinomial Distribution:} For a total of $n_{\text{shots}}$ trials, the set of counts $\{c_0, c_1, \dots, c_{2^{N+1}-1}\}$ for each of the $2^{N+1}$ basis states follows a multinomial distribution.
    \item \textbf{Central Limit Theorem (CLT):} The expectation value of an observable $\langle O \rangle$ is a sample mean calculated from a large number of shots. According to the CLT, for a sufficiently large $n_{\text{shots}}$, the sampling distribution of this mean is well-approximated by a Normal (Gaussian) distribution. This justifies the use of standard error as a measure of statistical uncertainty.
\end{itemize}

\subsection{Calculating Expectation Values and Statistical Uncertainty}

\subsubsection{Expectation Value}
An observable $O$ is measured by post-processing the counts from the simulation. The expectation value $\langle O \rangle$ is computed as:
\begin{equation}
    \langle O \rangle = \sum_{i=0}^{2^{N+1}-1} p_i \lambda_i
\end{equation}
where $p_i = c_i / n_{\text{shots}}$ is the estimated probability of measuring the basis state $|i\rangle$, and $\lambda_i = \langle i | O | i \rangle$ is the corresponding eigenvalue of the observable $O$.

\subsubsection{Statistical Uncertainty (Standard Error of the Mean)}
The statistical uncertainty of the estimated expectation value is given by the Standard Error of the Mean (SEM), denoted $\sigma_{\langle O \rangle}$. This is the quantity represented by the error bars in the plots. The SEM is calculated as:
\begin{equation}
    \sigma_{\langle O \rangle} = \frac{\sigma_O}{\sqrt{n_{\text{shots}}}}
\end{equation}
where $\sigma_O$ is the standard deviation of a \textit{single} measurement of the observable $O$. The variance, $\sigma_O^2$, is calculated directly from the simulation results using the formula:
\begin{equation}
    \begin{aligned}
        \sigma_O^2 &= \text{Var}(O) = \langle O^2 \rangle - \langle O \rangle^2 \\
        &= \left( \sum_{i} p_i \lambda_i^2 \right) - \left( \sum_{i} p_i \lambda_i \right)^2
    \end{aligned}
\end{equation}
This provides a direct method to quantify the statistical noise inherent in the Monte Carlo nature of the simulation.

\subsection{Application to Teleportation Observables}

We now apply this framework to the specific observables measured in the charge and energy teleportation protocols.

\subsubsection{Charge Teleportation (\texorpdfstring{$Q_B$}{QB})}
The charge at Bob's site is given by the operator $Q_B = \frac{1}{2}(I + Z_N)$. Since $Q_B$ is a projector, its eigenvalues are 0 and 1. The variance of a single measurement is:
\begin{equation}
    \text{Var}(Q_B) = \langle Q_B^2 \rangle - \langle Q_B \rangle^2 = \langle Q_B \rangle - \langle Q_B \rangle^2
\end{equation}
Using the relation $\langle Q_B \rangle = \frac{1}{2}(1 + \langle Z_N \rangle)$, we find:
\begin{equation}
    \begin{aligned}
        \text{Var}(Q_B) &= \frac{1}{2}(1 + \langle Z_N \rangle) - \frac{1}{4}(1 + \langle Z_N \rangle)^2 \\
        &= \frac{1}{4}(1 - \langle Z_N \rangle^2)
    \end{aligned}
\end{equation}
Thus, the SEM for the charge measurement is:
\begin{equation}
    \sigma_{\langle Q_B \rangle} = \frac{1}{2\sqrt{n_{\text{shots}}}} \sqrt{1 - \langle Z_N \rangle^2}
\end{equation}
Since this calculation relies on a single circuit execution (measuring in the Z-basis), the statistical analysis is straightforward.

\subsubsection{Energy Teleportation (\texorpdfstring{$H_B$}{HB})}
Bob's local Hamiltonian, for instance in the nearest-neighbor model ($H^{(2)}$), is $H_B = hZ_N + JX_{N-1}X_N$. The two terms, $Z_N$ and $X_{N-1}X_N$, do not commute and must be measured in separate circuit executions. This necessitates the use of error propagation to find the total uncertainty.

The expectation value is $\langle H_B \rangle = h\langle Z_N \rangle + J\langle X_{N-1}X_N \rangle$. Since the measurements for $\langle Z_N \rangle$ and $\langle X_{N-1}X_N \rangle$ are from independent sets of shots, their statistical errors are uncorrelated. The total variance of the mean is the sum of the individual variances:
\begin{equation}
    \begin{aligned}
        \sigma_{\langle H_B \rangle}^2 &= \text{Var}(h\langle Z_N \rangle) + \text{Var}(J\langle X_{N-1}X_N \rangle) \\
        &= h^2 \sigma_{\langle Z_N \rangle}^2 + J^2 \sigma_{\langle X_{N-1}X_N \rangle}^2
    \end{aligned}
\end{equation}
The individual SEMs are calculated as:
\begin{align}
    \sigma_{\langle Z_N \rangle}^2 &= \frac{1}{n_{\text{shots,Z}}} \left( \langle Z_N^2 \rangle - \langle Z_N \rangle^2 \right) \nonumber\\
    &= \frac{1}{n_{\text{shots,Z}}} \left( 1 - \langle Z_N \rangle^2 \right) \\
    \sigma_{\langle X_{N-1}X_N \rangle}^2 &= \frac{1}{n_{\text{shots,X}}} \left( 1 - \langle X_{N-1}X_N \rangle^2 \right)
\end{align}
The final expression for the uncertainty in the energy measurement is:
\begin{equation}
    \sigma_{\langle H_B \rangle} = \sqrt{
        \begin{aligned}    
        &\frac{h^2}{n_{\text{shots,Z}}} (1 - \langle Z_N \rangle^2) \\
        &+ \frac{J^2}{n_{\text{shots,X}}} (1 - \langle X_{N-1}X_N \rangle^2)
        \end{aligned}
    }
\end{equation}
This explicitly shows that the statistical uncertainty for the energy is a combination of noise from two independent experiments, quantitatively explaining the increased fluctuations observed in the Qiskit energy plots compared to the charge plots.

\subsection{Summary and Conclusions}

This analysis provides a rigorous basis for the statistical error bars and the discussion of noise in the main text. Key conclusions include:
\begin{itemize}
    \item The increased statistical noise in energy measurements is a direct consequence of error propagation from measuring non-commuting observables in separate experiments.
    \item The "meaningful statistical noise" observed for larger systems (e.g., $N=3, 4$ in the nearest-neighbor model) is a result of a decreasing Signal-to-Noise Ratio (SNR). As $N$ increases, the teleported signal $|\langle O \rangle|$ diminishes due to weaker correlations, while the single-shot variance $\sigma_O^2$ remains approximately constant (approaching 1). Consequently, the SEM $\sigma_{\langle O \rangle}$ becomes comparable to or larger than the signal itself, requiring a significantly larger $n_{\text{shots}}$ to resolve.
    \item The stability of the charge teleportation protocol is attributed to its reliance on a single measurement basis and a consistently strong signal across different coupling regimes.
\end{itemize}
